\documentclass[12pt]{article}
\usepackage[utf8]{inputenc}
\usepackage[T1]{fontenc}
\usepackage[margin=1in]{geometry}
\usepackage{setspace}
\usepackage{xcolor}
\usepackage[colorlinks=true, urlcolor=blue!60!black, linkcolor=blue!60!black]{hyperref}
\setlength{\parindent}{2em}
\setlength{\parskip}{1.25em}
\renewcommand{\baselinestretch}{1.0}
\begin{document}
\begin{center}
{\LARGE \textit{Griswold v. Connecticut}, 381 U.S. 479 (1965)\par}
\vspace{0.8em}
{\large Case Brief\par}
\end{center}
\vspace{1em}

\section{Memory Jogger}
In \href{https://www.courtlistener.com/opinion/107082/griswold-v-connecticut/}{\textit{Griswold v. Connecticut}}, the Supreme Court invalidated Connecticut's criminal ban on contraceptive use by married couples, recognizing a constitutional right to marital privacy derived from ``penumbras, formed by emanations'' of specific guarantees in the Bill of Rights.

\section{Detailed Case Facts}
Since 1879, Connecticut had maintained an anti-contraception law codified in \href{https://supreme.justia.com/cases/federal/us/381/479/}{Conn. Gen. Stat. \S~53-32 (1958 rev.)}, which provided that ``[a]ny person who uses any drug, medicinal article or instrument for the purpose of preventing conception shall be fined not less than fifty dollars or imprisoned not less than sixty days nor more than one year or be both fined and imprisoned.'' Under Connecticut's general accessory statute, \href{https://supreme.justia.com/cases/federal/us/381/479/}{Conn. Gen. Stat. \S~54-196 (1958 rev.)}, ``[a]ny person who assists, abets, counsels, causes, hires or commands another to commit any offense may be prosecuted and punished as if he were the principal offender.''

Prior judicial challenges to the statute had failed on threshold justiciability grounds. In \textit{Tileston v. Ullman}, 318 U.S. 44 (1943), a physician's declaratory-judgment suit was dismissed for lack of third-party standing, and in \textit{Poe v. Ullman}, 367 U.S. 497 (1961), the Supreme Court dismissed another declaratory challenge as unripe because the state had rarely enforced the statute (\href{https://supreme.justia.com/cases/federal/us/381/479/}{Justia Opinion}).

To force a direct constitutional test after \textit{Poe}, Estelle Griswold, Executive Director of the Planned Parenthood League of Connecticut, and Dr. C. Lee Buxton, a licensed physician, Yale Medical School professor, and Medical Director of the League, opened a birth control clinic in New Haven, Connecticut (\href{https://www.oyez.org/cases/1964/496}{Oyez}). The clinic operated from November 1 to November 10, 1961, during which Griswold and Buxton gave information, instruction, and medical advice to married persons on preventing conception, examined wives, and prescribed appropriate contraceptive devices or materials, typically charging a fee though serving some couples for free (\href{https://supreme.justia.com/cases/federal/us/381/479/}{Justia Opinion}). On November 10, 1961, Griswold and Buxton were arrested and charged as accessories to the unlawful use of contraceptives.

\section{Procedural History}
\begin{itemize}
    \item \textbf{Circuit Court for the Sixth Circuit of Connecticut (Jan. 2, 1962):} Following a bench trial, Griswold and Buxton were found guilty as accessories under Conn. Gen. Stat. \S\S~53-32 and 54-196 and fined \$100 each, over their defense that the accessory statute as applied violated the Fourteenth Amendment (\href{https://supreme.justia.com/cases/federal/us/381/479/}{Justia Opinion}; \href{https://en.wikipedia.org/wiki/Griswold_v._Connecticut}{Wikipedia Case Summary}).
    \item \textbf{Appellate Division of the Circuit Court (Jan. 7, 1963):} Affirmed the convictions (\href{https://supreme.justia.com/cases/federal/us/381/479/}{Justia Opinion}).
    \item \textbf{Supreme Court of Errors of Connecticut (1964):} Affirmed the convictions, \textit{State v. Griswold}, 151 Conn. 544, 200 A.2d 479 (1964) (\href{https://supreme.justia.com/cases/federal/us/381/479/}{Justia Opinion}).
    \item \textbf{Supreme Court of the United States (1965):} Noted probable jurisdiction, 379 U.S. 926 (1964); heard oral argument on March 29--30, 1965; and reversed the judgment on June 7, 1965, in \href{https://www.courtlistener.com/opinion/107082/griswold-v-connecticut/}{381 U.S. 479 (1965)} (\href{https://www.oyez.org/cases/1964/496}{Oyez}).
\end{itemize}

\section{Judicial Votes}
The Supreme Court reversed the convictions by a \textbf{7--2} vote (\href{https://supreme.justia.com/cases/federal/us/381/479/}{Justia Opinion}; \href{https://www.oyez.org/cases/1964/496}{Oyez}):
\begin{itemize}
    \item \textbf{Opinion of the Court (5 Justices):} Justice William O. Douglas (author), joined by Chief Justice Earl Warren, Justice Tom C. Clark, Justice William J. Brennan, Jr., and Justice Arthur J. Goldberg.
    \item \textbf{Concurring Opinion (joining the Opinion and Judgment of the Court):} Justice Arthur J. Goldberg (author), joined by Chief Justice Earl Warren and Justice William J. Brennan, Jr.
    \item \textbf{Concurrences in the Judgment Only (2 Justices):}
    \begin{itemize}
        \item Justice John Marshall Harlan II (author, concurring in the judgment).
        \item Justice Byron R. White (author, concurring in the judgment).
    \end{itemize}
    \item \textbf{Dissents (2 Justices):}
    \begin{itemize}
        \item Justice Hugo L. Black (author), joined by Justice Potter Stewart.
        \item Justice Potter Stewart (author), joined by Justice Hugo L. Black.
    \end{itemize}
\end{itemize}

\section{Holding}
\textbf{Judgment:} Reversed. The convictions of Griswold and Buxton were overturned, and Connecticut's statutory ban on the use of contraceptives by married couples was held unconstitutional (\href{https://www.courtlistener.com/opinion/107082/griswold-v-connecticut/}{CourtListener}).

\textbf{Rule Announced:} First, criminal defendants convicted as accessories for providing medical advice and contraceptives to married couples have third-party standing to assert the constitutional rights of those married persons (\href{https://supreme.justia.com/cases/federal/us/381/479/}{381 U.S. at 481}). Second, a state statute forbidding the use of contraceptives by married couples violates a fundamental constitutional right of marital privacy protected within the ``penumbras'' of specific guarantees of the Bill of Rights, made applicable to the states through the Fourteenth Amendment (\href{https://supreme.justia.com/cases/federal/us/381/479/}{381 U.S. at 484--486}).

\textbf{Holding vs. Dicta:}
\begin{itemize}
    \item \textbf{Holding:} (1) Accessories criminally convicted for assisting married couples have standing to raise the couples' constitutional privacy rights; and (2) states may not criminalize the \textit{use} of contraceptives by \textit{married couples}, because doing so sweeps unnecessarily broadly and directly invades the protected zone of marital privacy (\href{https://supreme.justia.com/cases/federal/us/381/479/}{381 U.S. at 481, 485--486}).
    \item \textbf{Dicta:} Justice Douglas's broad rhetoric describing marriage as ``intimate to the degree of being sacred'' and ``older than the Bill of Rights'' (\href{https://supreme.justia.com/cases/federal/us/381/479/}{381 U.S. at 486}), as well as his suggestion that a state might permissibly regulate the ``manufacture or sale'' of contraceptives rather than their use (\href{https://supreme.justia.com/cases/federal/us/381/479/}{381 U.S. at 485}), were not strictly necessary to decide the validity of Connecticut's anti-use statute. Likewise, Justice Goldberg's assurances that the decision ``in no way interferes with a State's proper regulation of sexual promiscuity or misconduct'' such as adultery and fornication (\href{https://supreme.justia.com/cases/federal/us/381/479/}{381 U.S. at 498--499}) were non-binding dicta later overtaken by decisions such as \textit{Eisenstadt v. Baird}, 405 U.S. 438 (1972), and \textit{Lawrence v. Texas}, 539 U.S. 558 (2003).
\end{itemize}

\section{Analysis of Opinions}

\subsection{Majority Opinion (Justice Douglas)}
Writing for five Justices, Justice Douglas first resolved standing by distinguishing \textit{Tileston v. Ullman}: unlike a physician seeking an abstract declaratory judgment, Griswold and Buxton suffered actual criminal convictions as accessories, giving them standing to argue that the underlying act they assisted could not constitutionally be a crime (\href{https://supreme.justia.com/cases/federal/us/381/479/}{381 U.S. at 481}).

On the merits, Douglas expressly declined to revive \textit{Lochner v. New York}, 198 U.S. 45 (1905), emphasizing that the Court does not sit as a ``super-legislature'' over economic or business regulations (\href{https://supreme.justia.com/cases/federal/us/381/479/}{381 U.S. at 482}). To avoid relying directly on unadorned substantive due process, Douglas reasoned that specific guarantees in the Bill of Rights carry peripheral rights---``penumbras, formed by emanations from those guarantees that help give them life and substance''---that create ``zones of privacy'' (\href{https://supreme.justia.com/cases/federal/us/381/479/}{381 U.S. at 484}). He drew on five amendments: the First Amendment (freedom of association and privacy in associations, as in \textit{NAACP v. Alabama}, 357 U.S. 449 (1958), and reinterpreted \textit{Meyer v. Nebraska} and \textit{Pierce v. Society of Sisters}), the Third Amendment (prohibition on quartering soldiers in any house), the Fourth Amendment (protection of persons, houses, papers, and effects against unreasonable searches and seizures), the Fifth Amendment (Self-Incrimination Clause creating a zone of privacy), and the Ninth Amendment (confirming that enumerated rights do not deny or disparage unenumerated rights retained by the people) (\href{https://supreme.justia.com/cases/federal/us/381/479/}{381 U.S. at 482--485}). Because Connecticut forbade the \textit{use} of contraceptives rather than merely regulating their manufacture or sale, it swept unnecessarily broadly and had a ``maximum destructive impact'' on the marital relationship, evoking the repulsive prospect of police searching ``the sacred precincts of marital bedrooms'' (\href{https://supreme.justia.com/cases/federal/us/381/479/}{381 U.S. at 485--486}).

\subsection{Concurring Opinion (Justice Goldberg, joined by Chief Justice Warren and Justice Brennan)}
Justice Goldberg joined the Court's opinion and judgment but wrote separately to highlight the Ninth Amendment. Relying on James Madison's statements during the First Congress and Justice Joseph Story's \textit{Commentaries}, Goldberg argued that while the Ninth Amendment is not an independent source of substantive rights, its text and history demonstrate the Framers' belief that fundamental personal liberties exist beyond those expressly enumerated in the first eight amendments (\href{https://supreme.justia.com/cases/federal/us/381/479/}{381 U.S. at 488--493}). To guard against unconstrained judicial discretion, judges must look to the ``traditions and [collective] conscience of our people'' (\href{https://supreme.justia.com/cases/federal/us/381/479/}{381 U.S. at 493}). Applying strict scrutiny, Goldberg concluded that Connecticut's asserted interest in discouraging extramarital relations was not compelling or narrowly tailored, because the state could directly punish adultery and fornication without intruding into marriage (\href{https://supreme.justia.com/cases/federal/us/381/479/}{381 U.S. at 497--498}).

\subsection{Concurrence in the Judgment (Justice Harlan)}
Justice Harlan declined to join Douglas's penumbral theory because he feared it implicitly tied Fourteenth Amendment Due Process to the Bill of Rights incorporation debate (\href{https://supreme.justia.com/cases/federal/us/381/479/}{381 U.S. at 499--500}). Incorporating his celebrated dissent in \textit{Poe v. Ullman}, 367 U.S. 497, 522 (1961), Harlan maintained that ``[t]he Due Process Clause of the Fourteenth Amendment stands, in my opinion, on its own bottom'' and directly protects basic values ``implicit in the concept of ordered liberty'' under \textit{Palko v. Connecticut}, 302 U.S. 319 (1937) (\href{https://supreme.justia.com/cases/federal/us/381/479/}{381 U.S. at 500}). Judicial restraint, Harlan argued, comes not from an artificial restriction to enumerated Bill of Rights text, but from respect for history, societal traditions, federalism, and separation of powers (\href{https://supreme.justia.com/cases/federal/us/381/479/}{381 U.S. at 501--502}).

\subsection{Concurrence in the Judgment (Justice White)}
Justice White also grounded his vote directly in Fourteenth Amendment substantive due process, citing \textit{Meyer}, \textit{Pierce}, \textit{Skinner v. Oklahoma}, and \textit{Prince v. Massachusetts} for the proposition that marital and family intimacies carry heightened protection (\href{https://supreme.justia.com/cases/federal/us/381/479/}{381 U.S. at 502--503}). Even accepting Connecticut's sole asserted justification---deterring illicit premarital and extramarital sexual relationships---White found the statute failed means-ends scrutiny: because Connecticut did not bar the possession or sale of contraceptives to prevent disease, banning married couples from using them bore no rational, let alone necessary, connection to deterring extramarital affairs (\href{https://supreme.justia.com/cases/federal/us/381/479/}{381 U.S. at 505--507}).

\subsection{Dissenting Opinion (Justice Black, joined by Justice Stewart)}
Justice Black stated that he found the Connecticut law personally offensive, yet insisted that no constitutional provision prohibited it (\href{https://supreme.justia.com/cases/federal/us/381/479/}{381 U.S. at 507--508}). First, he rejected any First Amendment defense because Griswold and Buxton engaged in conduct (running a clinic, examining patients, and dispensing devices), not pure speech (\href{https://supreme.justia.com/cases/federal/us/381/479/}{381 U.S. at 508}). Second, he attacked Douglas's ``right of privacy'' as an amorphous substitute that both dilutes and distorts specific constitutional text (\href{https://supreme.justia.com/cases/federal/us/381/479/}{381 U.S. at 509--510}). Third, he lambasted the Goldberg, Harlan, and White concurrences as a revival of \textit{Lochner}-era ``natural law due process,'' warning that empowering federal judges to strike down laws they deem unreasonable or contrary to the ``collective conscience'' turns the Court into a ``day-to-day constitutional convention'' and a ``bevy of Platonic Guardians'' (\href{https://supreme.justia.com/cases/federal/us/381/479/}{381 U.S. at 511--527}).

\subsection{Dissenting Opinion (Justice Stewart, joined by Justice Black)}
Justice Stewart famously opened by calling the Connecticut statute ``an uncommonly silly law,'' yet concluded that the Court had no authority to invalidate it (\href{https://supreme.justia.com/cases/federal/us/381/479/}{381 U.S. at 527}). Systematically walking through each of the six amendments cited by the majority, Stewart observed that no speech was abridged, no soldier was quartered, no search or seizure occurred, and no one was compelled to self-incriminate (\href{https://supreme.justia.com/cases/federal/us/381/479/}{381 U.S. at 528--529}). He dismissed Goldberg's reliance on the Ninth Amendment as ``turn[ing] somersaults with history,'' explaining that the Ninth and Tenth Amendments were designed to limit federal power, not empower federal courts to void state statutes (\href{https://supreme.justia.com/cases/federal/us/381/479/}{381 U.S. at 529--530}). The proper remedy, Stewart concluded, was political repeal by the Connecticut legislature (\href{https://supreme.justia.com/cases/federal/us/381/479/}{381 U.S. at 531}).

\section{Examples: Future Applications}
\begin{enumerate}
    \item \textbf{Same-Side (Invalid): State Ban on Contraceptive Distribution to Unmarried Adults.} A state enacts a felony statute barring the distribution of contraceptives to unmarried individuals while permitting distribution to married couples. Under \textit{Griswold}'s logic---and as the Court later held in \textit{Eisenstadt v. Baird}, 405 U.S. 438 (1972) (\href{https://en.wikipedia.org/wiki/Griswold_v._Connecticut}{Wikipedia})---the statute falls on the same side as Connecticut's ban because the core privacy interest in deciding whether to bear or beget a child cannot be arbitrarily restricted by marital status without running afoul of both privacy and equal protection principles.
    \item \textbf{Same-Side (Invalid): Compulsory State Sterilization or Two-Child Cap for Married Couples.} A state experiencing severe resource shortages passes a law requiring married couples to use state-mandated birth control or undergo sterilization after having two children. As Justice Goldberg explicitly noted in concurrence (\href{https://supreme.justia.com/cases/federal/us/381/479/}{381 U.S. at 496--497}), this directly invades the ``sacred precincts'' of marital intimacy and reproductive autonomy, failing strict scrutiny under \textit{Griswold}.
    \item \textbf{Opposite-Side (Valid): State Prescription and FDA Safety Regulations for Hormonal Contraceptives.} A state requires that oral hormonal contraceptives be dispensed only by licensed pharmacies pursuant to a physician's prescription after a safety screening. Unlike Connecticut's total ban on \textit{use}, this health-and-safety regulation targets commercial distribution and medical safety rather than outlawing marital intimacy, aligning with Justice Douglas's distinction between forbidding use and regulating manufacture or sale (\href{https://supreme.justia.com/cases/federal/us/381/479/}{381 U.S. at 485}).
    \item \textbf{Opposite-Side (Valid): Criminalizing Non-Consensual Intimate Conduct or Domestic Abuse Inside the Marital Home.} A spouse prosecuted for domestic battery or marital rape invokes \textit{Griswold}, arguing that the state cannot police conduct inside the marital bedroom. This falls squarely on the opposite side: \textit{Griswold}'s zone of privacy protects consensual, bilateral marital intimacy, not physical harm or coercion against a non-consenting partner, where the state has an indisputable compelling interest.
    \item \textbf{Fence-Sitter (Close Case): State Requirement of Parental Notification for Minors Obtaining Prescription Contraceptives.} A state allows minors aged 16 and 17 to obtain non-prescription contraceptives (cf. \textit{Carey v. Population Services International}, 431 U.S. 678 (1977)) but requires clinics to notify a parent before prescribing hormonal birth control to minors under 16. This sits on the fence because it pits the minor's evolving reproductive privacy interest derived from \textit{Griswold} against both the state's heightened interest in protecting minors' health and the parents' own traditional Fourteenth Amendment liberty interest under \textit{Meyer} and \textit{Pierce} in directing the upbringing of their children.
\end{enumerate}

\section{Critique}
\textbf{Originalist and Textualist Critiques:} Conservative and originalist scholars have long targeted \textit{Griswold}'s methodology as unmoored from constitutional text. In his influential 1971 article \href{https://scispace.com/pdf/neutral-principles-and-some-first-amendment-problems-euvfezrxrs.pdf}{\textit{Neutral Principles and Some First Amendment Problems}}, 47 Ind. L.J. 1 (1971), Robert H. Bork argued that Justice Douglas made an unjustified logical leap from specific Bill of Rights protections to a freestanding, unenumerated ``right of privacy,'' leaving courts with no principled way to distinguish protected ``sexual gratification'' from unprotected economic or moral choices. Similarly, in his concurring opinion in \textit{Dobbs v. Jackson Women's Health Organization}, 597 U.S. 215 (2022), Justice Clarence Thomas argued that the Court ``should reconsider'' all substantive due process precedents, explicitly naming \textit{Griswold} (\href{https://scholarship.richmond.edu/cgi/viewcontent.cgi?article=1607&context=pilr}{Richmond Public Interest Law Review}).

\textbf{Progressive, Feminist, and Process-Theory Critiques:} Even scholars sympathetic to \textit{Griswold}'s outcome have critiqued its doctrinal foundation. Process theorist John Hart Ely was famously skeptical of unenumerated substantive due process rights, preferring democracy-reinforcing judicial review grounded in representation and political participation---though modern scholars argue that reproductive autonomy itself enables equal democratic participation (\href{https://www.californialawreview.org/print/democracy-promoting-review}{California Law Review}). Feminist and gender-policy scholars have also noted the ``paradoxes of privacy'' in \textit{Griswold}: framing reproductive freedom around \textit{marital} privacy initially excluded unmarried women, risked shielding domestic subordination within the ``private'' sphere, and left reproductive rights doctrinally vulnerable to post-\textit{Dobbs} attacks (\href{https://genderpolicyreport.umn.edu/privacy-matters-the-legacy-of-griswold/}{University of Minnesota Gender Policy Report}). Doctrinally, later Courts largely abandoned Douglas's ``penumbras'' framework in favor of Justice Harlan's straightforward Fourteenth Amendment substantive due process approach.

\textbf{Competing Constitutional Values:} \textit{Griswold} embodies the enduring tension between \textit{democratic self-governance and judicial restraint} (requiring courts to defer to elected legislatures unless a clear constitutional text is violated) and \textit{individual autonomy and substantive liberty} (protecting intimate personal decisions from majoritarian overreach even when not enumerated word-for-word in the Constitution).

\section{Key Quotations}
\begin{quote}
``The foregoing cases suggest that specific guarantees in the Bill of Rights have penumbras, formed by emanations from those guarantees that help give them life and substance.''
\end{quote}
\noindent --- Justice William O. Douglas, Opinion of the Court, 381 U.S. at 484 (\href{https://supreme.justia.com/cases/federal/us/381/479/}{Justia}; \href{https://www.courtlistener.com/opinion/107082/griswold-v-connecticut/}{CourtListener}).

\begin{quote}
``Would we allow the police to search the sacred precincts of marital bedrooms for telltale signs of the use of contraceptives? The very idea is repulsive to the notions of privacy surrounding the marriage relationship.''
\end{quote}
\noindent --- Justice William O. Douglas, Opinion of the Court, 381 U.S. at 485--486 (\href{https://supreme.justia.com/cases/federal/us/381/479/}{Justia}; \href{https://www.courtlistener.com/opinion/107082/griswold-v-connecticut/}{CourtListener}).

\begin{quote}
``The Due Process Clause of the Fourteenth Amendment stands, in my opinion, on its own bottom.''
\end{quote}
\noindent --- Justice John Marshall Harlan II, concurring in the judgment, 381 U.S. at 500 (\href{https://supreme.justia.com/cases/federal/us/381/479/}{Justia}; \href{https://www.courtlistener.com/opinion/107082/griswold-v-connecticut/}{CourtListener}).

\begin{quote}
``I like my privacy as well as the next one, but I am nevertheless compelled to admit that government has a right to invade it unless prohibited by some specific constitutional provision.''
\end{quote}
\noindent --- Justice Hugo L. Black, dissenting, 381 U.S. at 510 (\href{https://supreme.justia.com/cases/federal/us/381/479/}{Justia}; \href{https://www.courtlistener.com/opinion/107082/griswold-v-connecticut/}{CourtListener}).

\begin{quote}
``I think this is an uncommonly silly law. ... But we are not asked in this case to say whether we think this law is unwise, or even asinine. We are asked to hold that it violates the United States Constitution. And that I cannot do.''
\end{quote}
\noindent --- Justice Potter Stewart, dissenting, 381 U.S. at 527 (\href{https://supreme.justia.com/cases/federal/us/381/479/}{Justia}; \href{https://www.courtlistener.com/opinion/107082/griswold-v-connecticut/}{CourtListener}).

\end{document}
